\documentclass[12pt,a4paper]{scrartcl}
\usepackage[utf8]{inputenc}
\usepackage[ngerman]{babel}
\usepackage{amsmath}
\usepackage{amsfonts}
\usepackage{amssymb}
\usepackage{makeidx}
\usepackage{graphicx}
\usepackage[pdftex]{hyperref}
\title{Bericht}
\subtitle{über die Implementation von Pocket Monster in JavaFX}
\author{von Ahmed, Amir, Matthias, Omar und Salim.}
%\subject{Pocket Monsters}
\titlehead{Programmieren II an der Leibniz Universität Hannover}
\date{16.~Juni -- 12.~Juli~2025}

\renewcommand*\thesection{\arabic{section}. Woche}

%-------------- Zusammenfassung
\newcommand{\zusammenfassung}[1]{{

\bigskip\noindent\hrulefill\quad\textsf{{\small \textit{Zusammenfassung}}}}\quad\hrulefill

\medskip\noindent #1 

\noindent\hrulefill}
%--------------


\hypersetup{
    pdftitle = {Bericht über Implementation von Pocket Monster},
    pdfauthor = {Ahmed Amine Kortas, Amir Abid, Matthias-Benjamin Ahrens, Omar Ennabli und Salim Nahdi},
    pdfsubject = {Bericht über die implementieren von Pocket Monsters in JavaFX im Rahmen der Gruppenaufgabe in Programmieren II},
    pdfkeywords = {Programmieren, JavaFX, Leibniz Universität Hannover, Gruppenaufgabe, Pocket Monsters},
    colorlinks = true
}


\begin{document}
\maketitle
%\tableofcontents

\section{26.06. bis 30.06.2025}
In der ersten Woche haben wir mit der Grundstruktur unseres JavaFX-Projekts begonnen. Ziel war es, ein spielbares Grundgerüst mit klarer Trennung zwischen Model, View und Controller zu erstellen. Die wichtigsten Aufgaben und Fortschritte waren:

\begin{description}
\item[Projektstrukturierung] Es wurde die MVC-Struktur und das Basislayout mit \texttt{MenuScene} und \texttt{PlayerSelectionScene} erstellt.
\item[Spielstart] Erste Menüs und die Möglichkeit zur Spielerauswahl wurden implementiert.
\item[Elementsystem \& Attack-Multiplikatoren] Das Element-System wurde hinzu gefügt, das den Schaden der Angriffe beeinflusst.
\item[Damage-Animation] Die Implementierung einer simplen Angriffsanimation zur Visualisierung von Attacken wurde begonnen.
\end{description}

\zusammenfassung{Die erste Woche diente dem Aufbau des Grundgerüsts mit Menüs, Spiellogik und ersten Animationen. Die Zusammenarbeit funktionierte gut, und das Projekt war früh lauffähig.}
\section{01.07. bis 05.07.2025}
In der zweiten Woche konzentrierten wir uns auf die Spielmechanik, Interaktion und visuelle Rückmeldungen. Die Hauptbeiträge:

\begin{description}
\item[Damage Animation V2] Verbesserte Übergänge und visuelle Effekte bei Angriffen.
\item[TextBox für Kampfstatus] Eine Textanzeige wurde hinzugefügt, die Kampfnachrichten wie „It’s super effective!“ anzeigt.
\item[Cooldown-System] Ultimative Angriffe haben jetzt Abklingzeiten, die korrekt verwaltet werden.
\item[Victory/Defeat-Screen] Am Ende eines Kampfes wird nun ein entsprechender Bildschirm mit Bild und Text eingeblendet.
\item[Spielerwechsel und Buttons] Buttons für Angriff, Spielerwechsel und Ultimate wurden eingeführt und mit der Logik verbunden.
\item[HealthBar-Update] Gesundheitsanzeigen wurden überarbeitet, um zwei Spieler je Seite korrekt darzustellen.
\end{description}

\zusammenfassung{Diese Woche war durch intensive Arbeit an Benutzerinteraktion und Spielfeedback geprägt. Das Spielgefühl wurde deutlich verbessert und viele Bugs wurden behoben.}

\newpage
\section*{Anleitung}

\end{document}